% Preparación autónoma de VII, 2026-09-30. Fragmento sin inputs adicionales.
% RESULTADOS_RECUPERADOS, pruebas reunidas sin modificación de los originales:
% /Users/ruben/Documents/excelencia academica/output/CUMPLIMIENTO_EDITORIAL_20260928/fuentes/VIII_ES/manuscrito/29_retorno_y_memoria_traslacional.tex
% /Users/ruben/Documents/excelencia academica/output/CUMPLIMIENTO_EDITORIAL_20260928/fuentes/VIII_ES/manuscrito/30_pantalla_y_respuesta.tex
% /Users/ruben/Documents/excelencia academica/output/CUMPLIMIENTO_EDITORIAL_20260928/fuentes/VIII_ES/manuscrito/30b_variacion_energia_memoria.tex
% /Users/ruben/Documents/excelencia academica/output/CUMPLIMIENTO_EDITORIAL_20260928/fuentes/VIII_ES/manuscrito/30c_composicion_corriente_conexion.tex
% /Users/ruben/Documents/excelencia academica/output/CUMPLIMIENTO_EDITORIAL_20260928/fuentes/VIII_ES/manuscrito/30d_realizacion_geometrica.tex
% /Users/ruben/Documents/excelencia academica/output/CUMPLIMIENTO_EDITORIAL_20260928/fuentes/VIII_ES/manuscrito/31_corriente_y_respuesta_cartan.tex
% /Users/ruben/Documents/excelencia academica/output/CUMPLIMIENTO_EDITORIAL_20260928/fuentes/VIII_ES/manuscrito/33_corriente_espinorial_y_densidad.tex
% Dependencias internas ya presentes: núcleo común, rutas, lectores y unidades.
% Las inversas de 31 están demostradas en 10, hmtcuatro:gea:torsion.
% Legendre, fuente completa y restricciones permanecen en 20.
% Se incluye de 33 sólo lo utilizado por 10+20, no sus estados cosmológicos.
\section{Fondements géométriques et matériels de l’action conjointe}
\label{hmtcuatro:ant:seccion}

Les graines et les opérations d’APP, l’orientation et le régime du TRIT et le transport du TPK produisent les routes enrichies sur lesquelles agissent les lecteurs du noyau commun. Feuillet, orientation, retenue, frontière et mémoire sont conservés conjointement ; le retour de phase ne réinitialise pas cet état. Les constructions du continu appartiennent au même objet et ne sont pas transformées ici en secteurs indépendants. La réalisation qui suit transporte vers la géométrie et la matière cette structure déjà construite. Ses unités et ses couplages sont des sorties HMT antérieures : aucune comparaison numérique n’intervient comme générateur.

Pour que l’action et sa réduction puissent être lues au sein de ce manuscrit, nous réunissons les opérations d’écran, la variation effective du transport et le courant spinoriel qu’elles utilisent. L’élimination de Cartan--Holst est démontrée dans \ref{hmtcuatro:gea:torsion} ; l’action totale et sa réduction canonique, dans \ref{hmtcuatro:accion-retroaccion-restricciones}. Les formules de cette section précisent les données et les calculs qui interviennent dans les deux.

\subsection{Incidence, écran et soudure}
\label{hmtcuatro:ant:pantalla}

La représentation cubique de la frontière d’une route possède les six faces \(F=\{(i,\epsilon):1\leq i\leq3,\ \epsilon=\pm1\}\). Sur \(H_F=\mathbb R^F\), l’opposition est \(Re_{i,\epsilon}=e_{i,-\epsilon}\). Avec
\[
 u=\sum_{i,\epsilon}e_{i,\epsilon},\qquad
 P_u=\frac{uu^{\mathsf T}}6,\qquad
 P_-=\frac{I-R}2,\qquad
 P_0=\frac{I+R}2-P_u ,
\]
les trois projecteurs sont orthogonaux, leur somme vaut \(I\) et leurs rangs sont \(1,3,2\). L’incidence entre faces et sommets \(\nu\in\{-1,1\}^3\) est \(M_{(i,\epsilon),\nu}=\mathbf1_{\{\nu_i=\epsilon\}}\).

\begin{proposition}[Écran de l’incidence orientée]
\label{hmtcuatro:ant:prop-pantalla}
On a
\[
 MM^{\mathsf T}=12P_u+4P_-,\qquad
 \ker M^{\mathsf T}=P_0H_F .
\]
Le quotient \(Q=H_F/P_0H_F\), identifié à \((P_u+P_-)H_F\), est de dimension quatre. La forme \(B=P_--P_u\) sur \(Q\) a pour signature \((-+++)\).
\end{proposition}
\begin{proof}
Chaque face contient quatre sommets ; deux faces opposées n’en partagent aucun et deux faces d’axes distincts en partagent deux. Par conséquent,
\[
 MM^{\mathsf T}=4I+2(uu^{\mathsf T}-I-R)
              =12P_u+4P_- .
\]
L’identité \(\|M^{\mathsf T}x\|^2=\langle x,MM^{\mathsf T}x\rangle\) détermine le noyau. Dans la base orthonormée
\[
 t=u/\sqrt6,\qquad d_i=(e_{i,+}-e_{i,-})/\sqrt2
\]
de \(Q\), la matrice de \(B\) est \(\operatorname{diag}(-1,1,1,1)\).
\end{proof}

Soit \(W_\square=\sqrt6P_u+\sqrt2P_-\). Alors
\[
 W_\square e_{i,\epsilon}=t+\epsilon d_i=:n_{i,\epsilon},
 \qquad B(n_{i,\epsilon},n_{i,\epsilon})=0,\qquad
 MM^{\mathsf T}|_Q=2W_\square^*W_\square|_Q .
\]
Ces égalités s’obtiennent en appliquant les projecteurs à la base des faces. Les rotations propres du cube préservent les projecteurs, l’orientation temporelle \(t\) et la forme \(B\).

La route enrichie fournit à chaque arête \(a\) sa face, son signe \(\sigma_m(a)\) et le transport \(U_m(a):Q_{s(a)}\to Q_{t(a)}\). Sa soudure cellulaire est
\begin{equation}
 \beta_m(a)=\ell_m\sigma_m(a)n_{\lambda(\underline a)}^{s(a)},
 \quad \ell_m=\ell_0\,9^{-m},\quad
 \beta_m(\bar a)=-U_m(a)\beta_m(a).
 \label{hmtcuatro:ant:soldadura}
\end{equation}
L’inversion \(U_m(\bar a)=U_m(a)^{-1}\) implique qu’inverser deux fois restitue la soudure. Le registre ordonné de la retenue est également conservé : il n’est pas remplacé par le vecteur d’arrivée.

\begin{proposition}[Naturalité de la soudure]
\label{hmtcuatro:ant:refinamiento-soldadura} Soit \(a=a_1\cdots a_9\) un raffinement compatible : les incidences filles transportées à l’origine du prédécesseur conservent face et signe, et \(\ell_{m+1}=\ell_m/9\). Si \(\mathcal U_j:Q_{m,s(a)}\to Q_{m+1,s(a_j)}\) inclut l’identification par raffinement et le transport jusqu’à \(s(a_j)\), alors
\[
 \beta_m(a)=\sum_{j=1}^9\mathcal U_j^{-1}\beta_{m+1}(a_j).
\]
\end{proposition}
\begin{proof}
Chaque terme dans la fibre d’origine est \((\ell_m/9)\sigma_m(a)n_{\lambda(\underline a)}\). Leur somme donne \eqref{hmtcuatro:ant:soldadura}. Inverser la route inverse l’ordre des transports et chaque transport ; la formule d’inversion de la soudure donne le signe global. La concaténation conserve dans le même ordre la mémoire des neuf incidences. L’égalité s’itère par composition de raffinements compatibles.
\end{proof}

\subsection{Représentation du transport et réalisation géométrique}
\label{hmtcuatro:ant:geometria}

Le registre d’une route \(\gamma\) compte les inversions et les changements effectifs de feuillet au moyen du cocycle \(c(\gamma)=(c_g(\gamma),c_s(\gamma))\in\mathbb Z^2\). Le décompte est cumulatif sur les routes vers l’avant ; son extension au groupoïde est orientée et attribue des incréments opposés à la route inverse. Pour des routes composables, \(c(\gamma_2\gamma_1)=c(\gamma_2)+c(\gamma_1)\), et \(c(\bar\gamma)=-c(\gamma)\). Soit \(R_4\) un élément d’ordre quatre dans \(\operatorname{Spin}^+(3,1)\) dont la représentation vectorielle \(\Lambda(R_4)\) fixe le vecteur temporel unitaire \(u_0\). La représentation affine reçue est
\[
 \rho_{\rm C}^{\rm disc}(\gamma)
 =\bigl(R_4^{c_g(\gamma)},\ell_0c_s(\gamma)u_0\bigr).
\]
Avec le produit \((R,a)(S,b)=(RS,a+\Lambda(R)b)\), la fixation de \(u_0\) et l’additivité de \(c\) prouvent directement que \(\rho_{\rm C}^{\rm disc}\) conserve l’identité, la concaténation et l’inversion. Lorsque la composante rotationnelle revient à l’identité, la translation de mémoire reste déterminée par \(c_s(\gamma)\), qui n’est pas réinitialisé.

La normalisation
\[
 N_{\ell_0}(R,a)=(\Lambda(R),a/\ell_0),\qquad
 \rho_{\rm C}=N_{\ell_0}\circ\rho_{\rm C}^{\rm disc}
\]
est également un homomorphisme, car \((a+\Lambda(R)b)/\ell_0=a/\ell_0+\Lambda(R)(b/\ell_0)\). Le relèvement \(R_4^{c_g(\gamma)}\) est conservé lorsque des spineurs agissent ; il ne se récupère pas à partir du seul \(\Lambda(R_4^{c_g(\gamma)})\), dont la représentation a pour noyau \(\{1,-1\}\).

Dans une réalisation d’écran \(J_{{\rm scr},m}\) qui transporte la forme et la soudure, \(e_m=J_{{\rm scr},m}\beta_m\). La réalisation lisse considérée consiste en une région orientée \(\mathcal U\) de dimension quatre, une cotétrade non dégénérée \(e\), une connexion de Lorentz \(\omega\) et une réalisation des routes \(\mathfrak R_{\rm C}\) avec la condition de compatibilité
\begin{equation}
 \operatorname{Hol}_{\mathcal A_{\ell_0}}
          (\mathfrak R_{\rm C}(\gamma))
 =\rho_{\rm C}(\operatorname{Hol}_{\rm TPK}(\gamma)),\qquad
 \mathcal A_{\ell_0}=
 \begin{pmatrix}\omega&\ell_0^{-1}e\\0&0\end{pmatrix}.
\label{hmtcuatro:ant:compatibilidad-holonomia}
\end{equation}
L’argument de \(\rho_{\rm C}\) conserve le cocycle de l’holonomie enrichie ; il n’est pas remplacé par sa seule phase réduite. L’échelle \(\ell_0>0\) est constante dans la coordinatisation. Ce sont les données et la condition de la réalisation lisse utilisée ; préserver la composition des chemins n’identifie pas entre eux des chemins distincts de mêmes extrémités et n’élimine pas leur mémoire. La métrique est \(g=\eta_{IJ}e^I\otimes e^J\), \(\eta=\operatorname{diag}(-1,1,1,1)\).

\begin{proposition}[Courbure, covariance et Bianchi]
\label{hmtcuatro:ant:curvatura} Avec \(F_\omega=\mathrm d\omega+\omega\wedge\omega\) et \(T=\mathrm de+\omega\wedge e\),
\[
 \mathcal F_{\ell_0}=
 \begin{pmatrix}F_\omega&\ell_0^{-1}T\\0&0\end{pmatrix},
 \qquad D_\omega T=F_\omega\wedge e,\quad D_\omega F_\omega=0 .
\]
Un changement de repère \(g_L:\mathcal U\to SO^+(3,1)\) agit par \(e'=g_Le\), \(\omega'=g_L\omega g_L^{-1}-\mathrm dg_L\,g_L^{-1}\) ; alors \(T'=g_LT\), \(F_{\omega'}=g_LF_\omega g_L^{-1}\).
\end{proposition}
\begin{proof}
La multiplication des matrices de formes donne le bloc diagonal \(\mathrm d\omega+\omega\wedge\omega\) et le bloc translationnel \(\ell_0^{-1}(\mathrm de+\omega\wedge e)\). En substituant les champs transformés, les termes contenant \(\mathrm dg_L\) s’annulent. Développer \(D_\omega T\) laisse \((\mathrm d\omega+\omega\wedge\omega)\wedge e\). Dans \(D_\omega F_\omega\), les dérivées de \(\omega\) et les produits cubiques s’annulent deux à deux.
\end{proof}

Le feuillet tritique \(\tau\in\{+1,0,-1\}\) conserve son régime dans la réalisation de Cartan avec les générateurs \(J_{ab},P_a^{(\tau)}\), l’action vectorielle de Lorentz et \([P_a^{(\tau)},P_b^{(\tau)}]=-\tau J_{ab}\). Pour \(\mathcal A_\tau=\frac12\omega^{ab}J_{ab}+\ell^{-1}e^aP_a^{(\tau)}\),
\[
 \mathcal F_\tau=
 \tfrac12(F_\omega^{ab}-\tau\ell^{-2}e^a\wedge e^b)J_{ab}
 +\ell^{-1}T^aP_a^{(\tau)} .
\]
En effet, le crochet de Lorentz donne \(F_\omega\), le crochet mixte donne \(D_\omega e\) et le crochet translationnel donne le terme \(-\tau e\wedge e/(2\ell^2)\). La condition \eqref{hmtcuatro:ant:compatibilidad-holonomia} correspond au feuillet central ; les feuillets extrêmes utilisent leur représentation \(\rho_{{\rm C},\tau}\) et leur condition d’holonomie propres.

\subsection{Différentielle du transport et courant matériel}
\label{hmtcuatro:ant:corriente-transporte}

Dans une réalisation cellulaire finie, soit \(U_a:\mathcal H_{s(a)}\to\mathcal H_{t(a)}\) un transport inversible entre fibres hermitiennes. On ne présuppose pas que tout transport de Lorentz soit unitaire pour le produit énergétique. On choisit une orientation par arête. Pour
\[
 (Df)_a=f_{t(a)}-U_af_{s(a)},\quad r=Df,\quad
 q=\mathsf W r,\quad v_a=U_af_{s(a)},\quad
 E=\langle r,\mathsf W r\rangle ,
\]
le poids \(\mathsf W=\mathsf W^*>0\) peut coupler des arêtes. Ce poids n’est pas \(W_\square\). Les variations correspondent à la représentation de la connexion et au problème de bord choisi.

\begin{proposition}[Variation complète]
\label{hmtcuatro:ant:variacion-completa} Pour une courbe différentiable de données, avec \(A_a=\dot U_aU_a^{-1}\), on a
\begin{equation}
 \dot E=2\operatorname{Re}\sum_a
 \langle q_a,\dot f_{t(a)}-U_a\dot f_{s(a)}-A_av_a\rangle
 +\langle r,\dot{\mathsf W}r\rangle .
 \label{hmtcuatro:ant:variacion}
\end{equation}
À champ et poids fixés, le gradient complet du transport est \(G_a=-2q_av_a^*\) :
\[
 \delta_U E=\sum_a\operatorname{Re}\operatorname{Tr}(G_a^*A_a).
\]
Dans la sous-classe unitaire \(A_a^*=-A_a\), sa partie antihermitienne \(J_a=v_aq_a^*-q_av_a^*\) suffit.
\end{proposition}
\begin{proof}
La règle du produit donne \(\dot E=2\operatorname{Re}\langle\mathsf W r,\dot r\rangle+
\langle r,\dot{\mathsf W}r\rangle\). Dériver chaque résidu produit \eqref{hmtcuatro:ant:variacion}. Comme \(\operatorname{Tr}(v_aq_a^*A_a)=q_a^*A_av_a\), on obtient \(G_a\). La partie hermitienne de \(G_a\) est orthogonale, pour \(\operatorname{Re}\operatorname{Tr}(X^*Y)\), aux variations antihermitiennes. Cette réduction ne s’applique que dans cette sous-classe.
\end{proof}

Pour une route avec \(U_\gamma=U_N\cdots U_1\), soit \(B_k=U_N\cdots U_{k+1}\). La règle du produit démontre
\begin{equation}
 \dot U_\gamma U_\gamma^{-1}
 =\sum_k B_kA_kB_k^{-1}.
 \label{hmtcuatro:ant:variacion-ruta}
\end{equation}
Par conséquent, un gradient \(G_\gamma\) est transporté vers l’arête \(k\) sous la forme \(G_k=B_k^*G_\gamma(B_k^{-1})^*\). La cyclicité de la trace prouve cette formule même si \(B_k\) n’est pas unitaire.

Si les coordonnées réelles de variation de la connexion sont \(\theta_b\), sa différentielle effective s’écrit \(A_a(\theta)=\sum_b B_{ab}\theta_b\). Pour \(S_{\rm mem}=\mathfrak a E\), à champ, poids et section d’action \(\mathfrak a\) fixés,
\begin{equation}
 \delta S_{\rm mem}=-\tfrac12\sum_b\theta_b s_b,\qquad
 s_b=4\mathfrak a\,\operatorname{Re}
          \sum_a\langle q_a,B_{ab}v_a\rangle .
 \label{hmtcuatro:ant:corriente-celular}
\end{equation}
Substituer \(A_a(\theta)\) dans la variation démontre l’égalité et fixe le facteur \(4\). Si l’appariement cellulaire orienté est une matrice inversible \(\mathsf M\), définie par \(\delta S_{\rm mem}=-\frac12\theta^{\mathsf T}\mathsf M\sigma\), le courant est \(\sigma=\mathsf M^{-1}s\). Cette matrice conserve les volumes et l’orientation ; elle ne peut être remplacée par une division scalaire sans fixer une base où celle-ci convient. Si \(f,\mathsf W,\mathfrak a\) varient également, tous les termes de \eqref{hmtcuatro:ant:variacion} et de \(E\delta\mathfrak a\) sont conservés.

Pour un poids par arête, l’inversion conserve la réponse complète en transformant conjointement
\[
 U_{\bar a}=U_a^{-1},\qquad
 r_{\bar a}=-U_a^{-1}r_a,\qquad
 \mathsf W_{\bar a}=U_a^*\mathsf W_aU_a .
\]
La substitution donne \(r_{\bar a}^*\mathsf W_{\bar a}r_{\bar a}=r_a^*\mathsf W_ar_a\). Sa dérivée conserve nécessairement \(\dot{\mathsf W}_{\bar a}
=U_a^*(A_a^*\mathsf W_a+\dot{\mathsf W}_a+\mathsf W_aA_a)U_a\) : éliminer ce terme altérerait le courant. Si le poids couple plusieurs arêtes, la même congruence est appliquée à la matrice complète avec le transport diagonal par blocs ; ses blocs croisés sont également conservés.

\begin{proposition}[Compatibilité variationnelle par raffinement]
\label{hmtcuatro:ant:refinamiento-corriente} Pour des applications fixées \(J_m,K_m\), supposons que, tout au long de la variation considérée,
\[
 D_{m+1}J_m=K_mD_m,\qquad
 K_m^*\mathsf W_{m+1}K_m=\mathsf W_m .
\]
Alors \(E_{m+1}(J_mf)=E_m(f)\) et leurs différentielles coïncident. Si les actions coïncident également et si \(P_m\) transporte les variations de connexion vers le raffinement, leurs courants vérifient
\[
 s_m=P_m^{\mathsf T}s_{m+1},\qquad
 \mathsf M_m\sigma_m=P_m^{\mathsf T}\mathsf M_{m+1}\sigma_{m+1}.
\]
\end{proposition}
\begin{proof}
Substituer \(D_{m+1}J_m=K_mD_m\) dans l’énergie, puis la congruence du poids, prouve la première égalité sur toute la courbe. La dériver et utiliser \(-\frac12\theta^{\mathsf T}s_m
=-\frac12(P_m\theta)^{\mathsf T}s_{m+1}\) prouve les autres. Si \(J_m\) varie, sa différentielle apporte \(2\operatorname{Re}\langle D_{m+1}J_mf,
\mathsf W_{m+1}D_{m+1}\dot J_m f\rangle\) ; elle n’est pas écartée.
\end{proof}

Le courant de mémoire ainsi obtenu provient de son action. Le courant spinoriel qui suit provient de l’action affine des champs. Ils se composent lorsqu’ils sont des termes d’une même action ; ils ne s’identifient pas du seul fait qu’ils partagent le transport. En particulier, éliminer une contorsion dont dépendent également le champ préparé, l’incidence ou son minimum exige de conserver ces différentielles. La formule linéaire de Cartan utilisée ci-après correspond au problème matériel affine expressément déclaré.

\subsection{Représentation spinorielle et courant affine}
\label{hmtcuatro:ant:espinores}

On fixe la structure de spin de la réalisation géométrique. Les matrices héritées, après complexification du module réel, sont
\begin{equation}
 \begin{gathered}
 J=\begin{pmatrix}0&-1\\1&0\end{pmatrix},\quad
 R=\begin{pmatrix}0&1\\1&0\end{pmatrix},\quad
 Z=\begin{pmatrix}1&0\\0&-1\end{pmatrix},\\
 g^0=J\otimes I_2,\quad g^1=R\otimes I_2,\quad
 g^2=Z\otimes R,\quad g^3=Z\otimes Z .
 \end{gathered}
 \label{hmtcuatro:ant:clifford}
\end{equation}
Leurs carrés sont \(-I_4,I_4,I_4,I_4\) et les matrices distinctes anticommutent. Il en résulte \(\{g^I,g^J\}=2\eta^{IJ}I_4\). Nous définissons
\[
 g_I=\eta_{IJ}g^J,\quad
 \mathbb B_{IJ}=\tfrac14[g_I,g_J],\quad
 A_D=-ig^0,\quad \bar\psi=\psi^\dagger A_D,\quad
 \Omega=\tfrac12\omega^{IJ}\mathbb B_{IJ}.
\]
Le développement du commutateur de trois matrices donne \([\mathbb B_{IJ},g^K]=\delta_J^Kg_I-\delta_I^Kg_J\). Par conséquent, \(\Omega\) représente la même connexion métrique, avec \(D\psi=\mathrm d\psi+\Omega\psi\) et \(D\bar\psi=\mathrm d\bar\psi-\bar\psi\Omega\). La matrice \(A_D\) est hermitienne et \(A_Dg^I\) est antihermitienne ; elles fixent l’adjoint et le facteur de l’action en signature \((-+++)\).

Soit \(\eta_I=\iota_{E_I}\operatorname{vol}_e\), de sorte que \(e^J\wedge\eta_I=\delta_I^J\operatorname{vol}_e\). L’action affine est
\begin{equation}
 S_D=\int_M\left\{
 \frac{\hbar}{2}
 [\bar\psi g^ID\psi-(D\bar\psi)g^I\psi]\wedge\eta_I
 -mc\,\bar\psi\psi\operatorname{vol}_e\right\}.
 \label{hmtcuatro:ant:accion-dirac}
\end{equation}
La masse et les sections positives \(c,\hbar\) sont les lectures reçues, et \([\psi]=L^{-3/2}\). La notation alternative \(\Gamma^I=-ig^I\), avec le terme principal \(i\hbar\Gamma^ID_I\), passe à la signature de Clifford opposée ; les deux conventions ne sont pas mélangées.

\begin{proposition}[Courant spinoriel normalisé]
\label{hmtcuatro:ant:corriente-spin} Avec \(B^{IJK}=\bar\psi g^{[I}g^Jg^{K]}\psi\), où l’antisymétrisation est normalisée, la convention \(\delta_\omega S_D=-\frac12\int\delta\omega^{JK}\wedge\sigma_{JK}\) donne
\begin{equation}
 \sigma_{JK}
 =-\frac{\hbar}{2}\bar\psi\{g^I,\mathbb B_{JK}\}\psi\,\eta_I
 =-\frac{\hbar}{2}B^I{}_{JK}\eta_I .
 \label{hmtcuatro:ant:sigma-spin}
\end{equation}
Le courant est indépendant de la contorsion lorsque \(e,\psi\) est fixé.
\end{proposition}
\begin{proof}
On a \(\delta\Omega=\frac12\delta\omega^{JK}\mathbb B_{JK}\). Les deux termes cinétiques de \eqref{hmtcuatro:ant:accion-dirac} donnent
\[
 \delta_\omega S_D=\frac{\hbar}{4}\int
 \delta\omega^{JK}\wedge\eta_I\,
          \bar\psi\{g^I,\mathbb B_{JK}\}\psi .
\]
La comparaison avec la convention variationnelle détermine le signe et le facteur. En développant l’anticommutateur à l’aide de Clifford, les contractions de degré un s’annulent et il reste \(\{g^I,\mathbb B_{JK}\}=g^{[I}g_Jg_{K]}\). L’action est linéaire en \(\Omega\), ce qui prouve l’indépendance.
\end{proof}

Si des connexions internes sont incluses, elles agissent sur le facteur matériel du module spinoriel, tandis que \(\mathbb B_{IJ}\) agit sur le facteur de spin. Ainsi, \([\,\mathbb B_{IJ}\otimes I,I\otimes\rho(A)\,]=0\). Faire varier \(\omega\) produit la même formule, contractée sur les composantes internes et sommée sur les multiplets présents. Cela préserve leurs représentations chirales, de couleur et de famille ; aucun multiplet n’est ajouté à ceux déjà construits.

\subsection{Contraction torsionnelle et termes croisés entre espèces}
\label{hmtcuatro:ant:contraccion}

L’inversion réelle de Cartan--Holst et de l’application \(T^I\mapsto T^I\wedge e^J-e^I\wedge T^J\), démontrée dans \ref{hmtcuatro:gea:torsion}, fixe sans ambiguïté la contorsion \(\mathfrak k_*\) pour le courant affine \eqref{hmtcuatro:ant:sigma-spin}. Ces formules conservent \(\gamma\in\mathbb R\setminus\{0\}\), une cotétrade non dégénérée et les sections \(G,c,\hbar,\gamma\) fixées pendant la variation. L’élimination utilise une seule fois l’interaction linéaire et donne \(-\frac14\mathfrak k_*^{IJ}\wedge\sigma_{IJ}\).

\begin{theorem}[Coefficient de la contraction spinorielle]
\label{hmtcuatro:ant:coeficiente} Soit \(\kappa=8\pi G/c^4\). L’interaction effective est
\[
 \mathcal L_{\rm spin}^{\rm ef}
 =C_\gamma q(B)\operatorname{vol}_e,\qquad
 C_\gamma=\frac{3\kappa c\hbar^2}{16}
                    \frac{\gamma^2}{1+\gamma^2},
\]
où
\[
 q(B)=\frac1{3!}B^{IJK}B_{IJK}
 =-(B^{012})^2-(B^{013})^2-(B^{023})^2+(B^{123})^2 .
\]
\end{theorem}
\begin{proof}
On calcule d’abord
\[
 \mathcal Y^{IJ}
 =\kappa c\frac{\gamma^2}{1+\gamma^2}
       (-\star+\gamma^{-1}I)\sigma^{IJ},\quad
 \mathcal B^I=\sum_J\iota_{E_J}\mathcal Y^{IJ},\quad
 T^I=\mathcal B^I-\tfrac14e^I\wedge\sum_L\iota_{E_L}\mathcal B^L,
\]
puis \(\mathfrak k_{IJK}=(T_{KIJ}-T_{IJK}-T_{JKI})/2\). Ce sont des compositions linéaires des inverses déjà démontrés. Pour expliciter tous les coefficients de la contraction, on ordonne le trivecteur sous la forme \((012,013,023,123)\) et l’on extrait le facteur \(\kappa c\hbar^2\gamma^2/(1+\gamma^2)\). La substitution de \(\sigma_{JK}/\hbar=-\frac12B^I{}_{JK}\eta_I\) dans ces compositions et dans \(-\frac14\mathfrak k^{IJ}\wedge\sigma_{IJ}\) donne
\[
 \begin{array}{c|rrrr|c}
  \text{opérateur sur }\sigma&012&013&023&123&
           \text{six coefficients de polarisation}\\ \hline
  -\star&-3/16&-3/16&-3/16&3/16&0\\
  I&0&0&0&0&0
 \end{array}
\]
Chaque entrée utilise \(e^I\wedge\eta_J=\delta_J^I\operatorname{vol}_e\) et les contractions écrites ci-dessus. Les quatre entrées diagonales et les six polarisations énumèrent la forme quadratique complète. La partie \(\gamma^{-1}I\) s’annule et la partie restante est \((3/16)q(B)\), ce qui prouve la formule.
\end{proof}

Pour plusieurs espèces, la variation s’effectue sur leur action commune : \(\sigma=\sum_s\sigma_s\). Comme la réponse \(\mathfrak k_*=\mathsf C_e(\sigma)\) est linéaire,
\begin{equation}
 -\tfrac14\mathsf C_e\!\left(\sum_s\sigma_s\right)
       \wedge\sum_t\sigma_t
 =-\tfrac14\sum_{s,t}\mathsf C_e(\sigma_s)\wedge\sigma_t .
 \label{hmtcuatro:ant:cruces-especies}
\end{equation}
Ce développement démontre que les termes \(s\ne t\) subsistent. Calculer des réponses séparées et ne conserver que \(s=t\) changerait l’action totale.

\subsection{Ordre normal du courant total}
\label{hmtcuatro:ant:car}

Dans la réalisation fermionique finie d’une cellule, on utilise \(\mathcal F=\bigoplus_n\Lambda^n V\), avec les relations \(\{a_r,a_s^\dagger\}=\delta_{rs}\), \(\{a_r,a_s\}=0\) et le vide d’occupation comme référence de l’ordre normal. Nous écrivons \(d\Gamma(M)=\sum_{r,s}a_r^\dagger M_{rs}a_s\) ; ce \(\Gamma\) désigne la seconde quantification, non l’holonomie nonadique. Pour le module spinoriel à quatre composantes, les matrices sont
\[
 \begin{array}{c|cccc}
 abc&012&013&023&123\\ \hline
 M_{abc}=A_Dg^ag^bg^c&
 iJ\otimes R&iJ\otimes Z&iI_2\otimes J&iZ\otimes J
 \end{array}.
\]
Multiplier \eqref{hmtcuatro:ant:clifford} donne ces matrices ; elles sont hermitiennes, de trace nulle et de carré \(I_4\).

\begin{proposition}[Contraction normale et conservation des termes croisés]
\label{hmtcuatro:ant:orden-normal} Pour toute matrice \(M\),
\[
 :d\Gamma(M)^2:
 =d\Gamma(M)^2-d\Gamma(M^2).
\]
Dans le module à quatre composantes,
\[
 \widehat{\mathscr Q}_{\rm sp}
 =\sum_{a=1}^4s_a d\Gamma(M_a)^2+2\widehat N,
 \qquad s=(-1,-1,-1,+1),\quad \widehat N=d\Gamma(I_4).
\]
L’opérateur est autoadjoint et conserve l’occupation. Dans une somme de multiplets, on emploie la même première identité avec les matrices totales \(\widetilde M_a\) ; les contractions entre espèces sont conservées.
\end{proposition}
\begin{proof}
Dans \(d\Gamma(M)^2\), la relation \(a_sa_t^\dagger=\delta_{st}-a_t^\dagger a_s\) sépare la contraction \(d\Gamma(M^2)\) du terme à deux créateurs et deux annihilateurs. Ce dernier est, par définition, le produit normal. Comme \(M_a^2=I_4\) et \(\sum_a s_a=-2\), la soustraction des contractions vaut \(+2\widehat N\). L’hermiticité des matrices rend les termes autoadjoints et chaque monôme conserve l’occupation.

Si \(\widetilde M=\bigoplus_s M_s\), alors \(d\Gamma(\widetilde M)=\sum_s d\Gamma(M_s)\) et \(d\Gamma(\widetilde M^2)=\sum_s d\Gamma(M_s^2)\). Les bilinéaires de multiplets distincts commutent ; en élevant la première somme au carré, il reste \(2d\Gamma(M_s)d\Gamma(M_t)\) pour \(s<t\). La soustraction d’une occupation n’élimine pas ces termes croisés. C’est la version opératorielle de \eqref{hmtcuatro:ant:cruces-especies}.
\end{proof}

L’invariance sous les changements unitaires internes s’obtient parce qu’ils agissent sur les indices matériels et entrelacent les matrices de courant : la seconde quantification transporte la formule complète par conjugaison. La densité et son évaluation sur un état utilisent en outre le volume de la cellule et l’état matériel déclarés. Le courant moyen et son second moment ne sont pas identifiés. Ces données complètent les fondements utilisés par l’action conjointe, sans introduire une seconde élimination de la torsion ni modifier ses conditions de bord.
